
In order to understand the motivations, limitations and behaviour of the experiment we need to know about a few key concepts in fluid dynamics.

%\subsection{Navier Stokes}


\subsection{Reynolds Number}
The Reynolds number ($\operatorname{Re}$) is a dimensionless number describing the ratio of inertial forces to viscous forces in a flow. For flow in a pipe it is defined by Pedley~\cite{introfluid} as


\begin{equation}\label{eq:reynolds}
\operatorname{Re} = \frac{U L \rho}{\mu}
\end{equation}
where $U$ is the characteristic velocity of the flow, $L$ is the characteristic length, $\rho$ is the density and $\mu$ is the dynamic viscosity. This is used to estimate the 'regime' of the flow, of which there are two primary types. 
\begin{enumerate}
\item The laminar regime, where viscous forces dominate over inertial forces
\item The turbulent regime where inertial forces dominate.
\end{enumerate}

The Reynolds number is the ratio between the inertial forces and viscous forces, so we get more turbulent flow for larger 
Reynolds numbers. A simple visual characterization of the flow types can be seen in figure \ref{fig:laminar_flow}. 
 Flow with $\operatorname{Re}\ll 1$ is referred to as \emph{Stokes flow} where we can ignore inertial forces completely. This has has the implication that not only is the flow completely guaranteed to be laminar, but it is time reversible. 
 
 If there are particles in the flow they will also revert perfectly as long as they too have a Reynolds number that is much smaller than 1~\cite{introfluid3}. 

\begin{figure}[H]
\centering
\includegraphics[width=0.6\textwidth]{figures/theory/laminarFlow.pdf}
\caption{This shows the principal difference between laminar and turbulent flow.}
\label{fig:laminar_flow}
\end{figure}


\subsection{Stokes Drag and Stokes' law}
The drag force $F_D$ exerted by a fluid on a spherical particle for $\operatorname{Re} \ll 1$ is found using  Stokes's law \cite{introfluid2}

\begin{equation}
F_D = 6\pi \mu R v,
\end{equation}

where $v$ is the velocity of the sphere relative to the fluid, $\mu$ is the dynamic viscosity and $R$ is the radius of the sphere. The terminal velocity of a sphere sinking in a liquid is found by equating the gravitational force $F_G$ acting on the sphere with the drag force. $F_G$ is calculated as

\begin{equation}
F_G = \Delta \rho g\cdot \frac{4\pi R^3}{3},
\end{equation}

where $\Delta \rho$ is the difference in density and $g$ is the specific gravity. We find that the terminal velocity of a sinking (or floating) sphere is

\begin{equation}\label{eq:fallingSphere}
v_s = \frac{2}{9} \frac{\Delta \rho}{\mu} g R^2.
\end{equation}
	%\subsection{Shear}
%
%\subsection{P\'{e}clet Number}
%The p\'{}clet number describes the ratio of thermal noise to other stuff. I really don't know anything about the piclet number.
